\newpage
\section{无标签中子电流档案的可追溯观测与证据约束解释}

本章说明核电厂运行环境中积累的中子电流历史记录与核仪表技术资料如何进入同一诊断辅助流程：前者用于形成可复核的信号现象，后者用于提供可定位的文献证据；本章进一步明确两者的任务分工及其受控连接方式。

\subsection{研究对象与数据结构}

测量资源来自同一核电厂采集周期的中子电流历史记录，共含来自四套仪表组的 42 条多通道 CSV 记录。每条记录包括时间列和 7 个中子电流通道；典型记录约含 $6\times10^{5}$ 个样本，采样间隔约为 1 s，并覆盖多日窗口。组内通道具有相对布置关系，因此其顺序不是任意的特征排列。原始导出文件可能含有局部时间戳紊乱或非单调片段，时间顺序校正因而是后续分析的前提。

\subsection{测量侧：从原始记录到观测包}

测量侧产出的是\textbf{观测包}，记录对哪条记录、以哪些确定性规则观察到哪些数据质量与信号现象。确定性脚本将结构化观测写入记事板，其中稳定字段形成供检索使用的紧凑告警字段（compact alert summary）；测量侧代理仅据此组织面向工程师的状况报告，不确认物理机理、设备状态或根因。

每次分析以一条记录为单位。初始化脚本解析时间列、按时间排序并生成唯一的活跃记录，同时记录样本数、名义采样间隔、缺失和不活跃通道等数据健康信息。后续分析只读取该记录，从而保留输入、参数、脚本版本和派生产物之间的关联；通道标识和组内相对关系仅用于组织观察，不被延伸为预设故障知识。

初始化后，脚本仅执行预先注册的分析工具，覆盖数据健康、形态与共变、事件与区段、同步与滞后、极值定位和截面快照；各工具将输入范围、参数与结构化输出写入观测包。图~\ref{fig:probe-examples} 展示其中三类统计现象；相关与滞后只描述共享变化、近同步或时间错位，不赋予仪表关系或因果方向。

为使事件识别可重复，当前实现对每个通道采用长度为 31 个样本的中心滑动窗口；记录边界处使用可获得的窗口样本。令 $W_t$ 为时刻 $t$ 的窗口，$x_{t,c}$ 为通道 $c$ 的电流值，则局部基线、点态绝对残差及其稳健局部尺度定义为

\[
\begin{aligned}
m_{t,c} &= \operatorname{median}_{\tau\in W_t}x_{\tau,c},\\
r_{t,c} &= \left|x_{t,c}-m_{t,c}\right|,\\
d_{t,c} &= \operatorname{median}_{\tau\in W_t}r_{\tau,c},\\
R_c &= \max_t x_{t,c}-\min_t x_{t,c}.
\end{aligned}
\]

这类以中位数和绝对偏差抑制局部异常影响的稳健估计思想，可追溯至相关基础研究 \citep{hampel1974influence,rousseeuw1993alternatives}。对于非恒定通道，突刺候选定义为同时满足局部与全程幅度条件的点：

\[
I^{\mathrm{spike}}_{t,c}=
\mathbb{I}\!\left(r_{t,c}>5d_{t,c}\right)
\mathbb{I}\!\left(r_{t,c}>0.05R_c\right).
\]

连续阳性点合并为短时偏离事件，长度不超过两个采样点的连续零值段另行记录；“突刺”仅指检测事件。两通道以 Pearson 相关系数描述共变，当前实现将 $r\geq0.99$ 标记为近冗余，将 $0.85\leq r<0.99$ 标记为强共变。这些阈值用于稳定生成 observation packet，而非故障判据。

各工具将输入范围、参数和结构化结果写入 observation packet 与记事板。紧凑告警字段直接抽取异常通道与计数、数据完整性、共变或滞后、极值和快照等稳定字段，用作检索输入；condition report 则按时间和通道上下文组织同一批已保存观察，供工程师复核。两者的区分兼顾了检索输入的可复现性与人工审查所需的上下文。

\subsection{证据侧：从结构化摘要到候选解释}

技术证据资源由核仪表与控制相关资料构成，并为每份文档保留来源 metadata，以支持段落定位和后续引用。当前资料库包含 64 份文档，其中 13 份构成以 IAEA 在线监测和仪表通道评估资料为主的权威核心，其余为 NRC、EPRI/DOE、OECD/NEA 及实验室或专题报告等同领域补充资料。在线监测与仪表通道评估的工程背景由 IAEA 指南和相关综述提供 \citep{iaea2008olm1,iaea2008olm2,hashemian2011online}。

核心资料与补充资料并非“相关”与“不相关”的二分：补充资料扩大技术覆盖范围，核心资料对应部署时预先指定的权威来源。证据侧既检索与观察相符的段落，也记录前 10 个结果中核心资料的比例，以检验是否满足来源策略。

紧凑告警字段是测量侧与证据侧之间唯一的受控连接：它将突刺、零值、共变或滞后映射为待核查的技术主题。查询措辞由作者依据摘要、语种和资料术语人工设定并复核。

\begin{table}[H]
\centering
\footnotesize
\setlength{\tabcolsep}{2pt}
\begin{tabular}{p{0.29\columnwidth} p{0.63\columnwidth}}
\toprule
摘要中的观察线索 & 检索意图与优先查找的证据 \\
\midrule
突刺或孤立零值 & 信号完整性和异常测量背景；定义、核查步骤与诊断注意事项 \\
共变或滞后 & 冗余、共享上下文与时序行为；仪表组织、在线监测与时间行为说明 \\
形态或数据健康线索 & 稳定性、噪声、可用性与监测实践；数据质量和监测指南 \\
极值或通道快照 & 运行上下文与通道间比较；运行指导和推荐检查项 \\
\bottomrule
\end{tabular}
\caption{由 observation packet 中的稳定字段形成受控检索意图。}
\label{tab:retrieval-intent}
\end{table}


该映射限制的是检索范围，而不是解释结论。例如，观察到某次零值事件时，系统可查找有关测量完整性或通道检查的资料段落；它不能据此自动声称出现了某一种仪表失效。

资料库中的 PDF 经必要裁切和文本清理后转换、分块并建立向量索引；每个 chunk 保留文档身份、来源类别和页码或段落定位。检索返回与观察现象相关的可定位段落，候选解释以 evidence-linked memo 组织，使每项表述关联观测字段、资料段落和引文。检索配置影响证据的可见性与引用方式，而非工程判断本身。
